\section{Implications and Limitations}\label{sec:implications}

\subsection{What would make ASML a card worth holding}

The model identifies which policies would raise Europe's leverage and which would not. Several cut against positions common in the European debate.

\emph{Reduce conditional exposure, not dependence as such.} Rearmament raises Europe's leverage over ASML only insofar as it lowers $\kappa S$, the part of Europe's security that Washington could withhold after realignment. Capabilities that Washington could not withdraw selectively, and an independent European deterrent above all, matter more for leverage than aggregate spending targets. By the same logic, a United States that values European security for its own sake, and would defend Europe whatever it did with ASML, would have less leverage over the chokepoint, not more.

\emph{Shrink the nesting, at a price.} Europe's exit threat is limited by US jurisdiction over ASML's subsystems, staff, and customers (\cref{prop:nested}a). Moving the EUV light source and US-origin components out of US jurisdiction would raise the exit premium. It would also be costly, slow, and read in Washington as preparation for exit, which is the trigger for a preventive demand (\cref{prop:control}).

\emph{Treat servicing as the lever.} Europe's one source of leverage that does not run through China is the power to interrupt servicing during a standoff (\cref{prop:holdup}). It works only if Europe can outlast the US buffer, which makes compute on European soil and stored model weights relevant. It also requires that Europe be able to commit to keeping servicing interrupted round by round, which in practice means a procedure that can order an interruption quickly and sustain it; EU procedures are slow where US instruments act immediately.

\emph{Expect a preventive demand if Europe's premium rises.} Any policy that succeeds in raising Europe's exit premium also raises Washington's incentive to demand control while its threats still bite. A Europe that anticipates this can bind itself to Atlantic alignment through an allied export-control compact that keeps ownership in European hands (\cref{rem:bind}). Doing so gives up the exit threat in exchange for keeping the asset. The choice between a credible promise and a credible threat is the central strategic decision the model puts to European governments.

\emph{De-risk from China's threats.} Reducing China's coercive reach over European supply chains never lowers Europe's payoff (\cref{prop:statics}d). Reducing the value of China's promises lowers it only when Europe already has leverage, which on our reading of \cref{fig:breakeven} it has not had.

For Washington the lessons are symmetric. Its leverage over ASML depends on a guarantee that is reliable for loyal allies, and every signal that it is not erodes that leverage. Rationing loyal allies' access to frontier AI destroys the AI lever. Firm-level penalties against ASML are a weak threat because carrying them out would damage America's own supply. A preventive grab for control requires a public commitment and risks a standoff, while an allied compact would secure the same alignment at lower cost.

\subsection{Limitations}

The limitation that should most worry a hostile reader is measurement. Every result is stated in terms of quantities, such as conditionality, China's credible offer, and the buffers in a standoff, that are hard to observe, and a negative exit premium can rationalize compliance after the fact. We have tried to answer this with predictions about instruments, timing, and pivots rather than outcomes, and with the order-of-magnitude comparison in \cref{fig:breakeven}. A sharper test would measure how European elites assess the conditionality of the US guarantee, for instance in a survey experiment that varies whether US protection is described as contingent on export-control alignment, and relate those assessments to support for resisting US export demands.

Three other choices limit the analysis. The great powers are unitary. Congress, the courts, and US firms that sell to China shape the credibility of American threats, and the MATCH Act shows members of Congress trying to act as a commitment device. Most results assume complete information. Uncertainty about Washington's resolve, about Europe's sovereignty costs, or about China's progress would produce standoffs more often than the baseline predicts, as \cref{prop:breakdown} illustrates. Finally, the stakes of the AI race are exogenous. If frontier AI conferred a decisive and lasting advantage, the values of the chokepoint to both great powers would rise sharply and the preventive logic of \cref{sec:control} would apply with more force. The model does not say whether that premise is correct.
